% Generated draft: requires mathematical and source review.
\documentclass[11pt]{article}
\usepackage[margin=25mm]{geometry}
\usepackage{fontspec}
\setmainfont{Latin Modern Roman}
\usepackage{amsmath,amssymb,mathtools,unicode-math}
\setmathfont{Latin Modern Math}
\usepackage{enumitem,needspace,xurl,hyperref}
\setlength{\parindent}{0pt}
\setlength{\parskip}{5pt}
\setlength{\emergencystretch}{3em}
\newcommand{\N}{\mathbb{N}}
\newcommand{\Z}{\mathbb{Z}}
\newcommand{\Q}{\mathbb{Q}}
\newcommand{\R}{\mathbb{R}}
\newcommand{\C}{\mathbb{C}}
\newcommand{\F}{\mathbb{F}}
\newcommand{\A}{\mathbb{A}}
\newcommand{\PP}{\mathbb{P}}
\newcommand{\E}{\mathbb{E}}
\newcommand{\eps}{\varepsilon}
\newcommand{\dd}{\,\mathrm{d}}
\begin{document}
\section*{UnsolvedMath ID: 30000207}
\textbf{Source problem code:} catalogue entry
\subsection{Definitions and mathematical statement}
Fix the Maass forms $\phi$ and $\phi'$ occurring in the source, and let $(\phi_i)$ be the intended family of Maass forms on the same space $Y$, with spectral parameters $\lambda_i$ and measure $\dd v$. Define
\[
 c_i=\int_Y\phi\phi'\phi_i\,\dd v,
 \qquad d_i=\gamma(\lambda_i)|c_i|^2,
\]
where $\gamma$ denotes the source's standard gamma-factor normalization. The integral is reproduced without inserting complex conjugation. The conjectural bound is
\[
 \forall\varepsilon>0\quad d_i\ll_{\phi,\phi',Y,\varepsilon}|\lambda_i|^\varepsilon
 \qquad (|\lambda_i|\longrightarrow\infty).
\]
Thus, for each $\varepsilon>0$, with the background data fixed, there should exist $C_\varepsilon>0$ and $T_\varepsilon>0$ such that $d_i\le C_\varepsilon|\lambda_i|^\varepsilon$ for every index $i$ with $|\lambda_i|\ge T_\varepsilon$. This expresses the usual asymptotic reading of $\ll$; it does not impose a bound at a possibly zero spectral parameter.

The supplied record does not specify $Y$, the automorphy group, the measure normalization, the normalization and completeness of $(\phi_i)$, the relation between $\lambda_i$ and the Laplace eigenvalue, or the formula for $\gamma$. These are essential unresolved parts of the formulation. In particular, an $L^2$ normalization or an explicit gamma factor must be recovered from the report rather than invented here.

The original excerpt places this conjecture after a bound with exponent $5/3$, comparing it with a convexity exponent $2$ for a degree-eight triple $L$-function. It describes the method as explicit principal-series calculations in a circle model followed by stationary phase, refers to an unidentified reference [1], and discusses difficulties for higher rank and discrete series. Its printed group is $PGL(2,\Z)$; this potentially corrupted group specification is not silently replaced by a real Lie group. The excerpt's historical claims about the first degree-eight subconvexity bound and earlier results of degree at most four are retained as source context, not independently established assertions.
\subsection{Short English statement}
For fixed Maass forms, the conjecture asks whether their gamma-normalized squared triple periods grow more slowly than every positive power of the varying spectral parameter. The restored bound is precise only after the geometric, spectral, and gamma-factor conventions have been recovered from the source.
\subsection{Sources}
\begin{itemize}
\item Local UnsolvedMath registry, permanent ID 30000207, source code catalogue entry, title ``Subconvexity Coefficient Bounds for Triple $L$-Functions.'' No proposer is supplied; proposed year 2005.
\item \emph{Enveloping Algebras and Geometric Representation Theory} (2005), Oberwolfach Reports, pp.\ 723--777, DOI 10.4171/owr/2005/13. Supplied source and literature reference: \url{https://doi.org/10.4171/owr/2005/13}.
\item The original excerpt cites [1], but the supplied record gives no bibliographic identity for it.
\item Supplied statement status: corrected\_verified. The verification note says the official report was checked directly and the triple-period coefficients and normalized quantities were restored. This is a supplied verification claim, not a check performed for this draft.
\item Supplied literature triage, checked 2026-08-21: open. Its assessment says the extracted $d_i$ bound is truncated and does not state a complete conjecture; that assessment does not account explicitly for the later restored statement. The record was created 2026-08-25 and updated 2026-09-16. Time-sensitive status and cited claims require independent verification.
\item Registry curation: Ulam AI / UnsolvedMath Contributors, CC BY 4.0; the underlying report retains its supplied provenance. This is an unreviewed statement-only draft.
\end{itemize}
\end{document}
